\documentclass[ug.tex]


\begin{document}

\chapter{Electromagnetic Theory}

\boxx{
\textbf{Maxwell Equations:}
\begin{itemize}
    \item Derivation of Maxwell’s Equations.
    \item Displacement Current.
    \item Boundary Conditions at Interface between Different Media.
    \item Wave Equations.
    \item Plane Waves in Dielectric Media.
    \item Poynting Theorem and Poynting Vector. (12 Lectures)
\end{itemize}

\textbf{EM Wave Propagation in Unbounded Media:}
\begin{itemize}
    \item Plane EM Waves Through Vacuum and Isotropic Dielectric Medium.
    \item Transverse Nature of Plane EM Waves.
    \item Refractive Index and Dielectric Constant.
    \item Wave Impedance.
    \item Propagation Through Conducting Media.
    \item Relaxation Time.
    \item Skin Depth. (12 Lectures)
\end{itemize}

\textbf{EM Wave in Bounded Media:}
\begin{itemize}
    \item Boundary Conditions at a Plane Interface Between Two Media.
    \item Reflection \& Refraction of Plane Waves at Plane Interface Between Two Dielectric Media.
    \item Laws of Reflection \& Refraction.
    \item Fresnel's Formulae for Perpendicular \& Parallel Polarization Cases.
    \item Brewster's Law. (12 Lectures)
\end{itemize}
}
\boxx{
\textbf{Polarization of Electromagnetic Waves:}
\begin{itemize}
    \item Description of Linear, Circular, and Elliptical Polarization.
    \item Uniaxial and Biaxial Crystals.
    \item Double Refraction.
    \item Polarization by Double Refraction.
    \item Nicol Prism.
    \item Ordinary \& Extraordinary Refractive Indices.
    \item Production \& Detection of Plane, Circularly, and Elliptically Polarized Light.
    \item Phase Retardation Plates: Quarter-Wave and Half-Wave Plates.
    \item Babinet Compensator and its Uses.
    \item Analysis of Polarized Light. (16 Lectures)
\end{itemize}

\textbf{Rotatory Polarization:}
\begin{itemize}
    \item Optical Rotation.
    \item Biot’s Laws for Rotatory Polarization.
    \item Fresnel’s Theory of Optical Rotation.
    \item Calculation of Angle of Rotation.
    \item Experimental Verification of Fresnel’s Theory.
    \item Specific Rotation. (8 Lectures)
\end{itemize}

}


\subsection*{I. Derivation of Maxwell’s Equations}

\subsubsection*{Statements:}

\begin{enumerate}
    \item \textbf{Introduction to Maxwell’s Equations:}
    \begin{itemize}
        \item Provide an overview of Maxwell's equations and their significance in describing electromagnetic phenomena.
    \end{itemize}
    
    \item \textbf{Gauss's Law for Electricity:}
    \begin{itemize}
        \item Present Gauss's law for electricity, emphasizing the relationship between electric charge and electric flux.
    \end{itemize}
    
    \item \textbf{Gauss's Law for Magnetism:}
    \begin{itemize}
        \item Present Gauss's law for magnetism, highlighting the absence of magnetic monopoles.
    \end{itemize}
    
    \item \textbf{Faraday's Law of Electromagnetic Induction:}
    \begin{itemize}
        \item Introduce Faraday's law, describing the electromotive force induced in a closed circuit by a changing magnetic field.
    \end{itemize}
    
    \item \textbf{Ampère's Law with Maxwell’s Addition:}
    \begin{itemize}
        \item Present Ampère's law and introduce Maxwell's addition—displacement current—to complete the set of Maxwell’s equations.
    \end{itemize}
\end{enumerate}

\subsection*{II. Displacement Current}

\subsubsection*{Statements:}

\begin{enumerate}
    \item \textbf{Introduction to Displacement Current:}
    \begin{itemize}
        \item Define displacement current as a concept introduced by Maxwell to account for changing electric fields in Ampère's law.
    \end{itemize}
\end{enumerate}

\subsection*{III. Boundary Conditions at Interface Between Different Media}

\subsubsection*{Statements:}

\begin{enumerate}
    \item \textbf{Boundary Conditions for Electric Field:}
    \begin{itemize}
        \item Outline the conditions governing the behavior of electric fields at the interface between different media.
    \end{itemize}
    
    \item \textbf{Boundary Conditions for Magnetic Field:}
    \begin{itemize}
        \item Outline the conditions governing the behavior of magnetic fields at the interface between different media.
    \end{itemize}
\end{enumerate}

\subsection*{IV. Wave Equations}

\subsubsection*{Statements:}

\begin{enumerate}
    \item \textbf{Wave Equations for Electric and Magnetic Fields:}
    \begin{itemize}
        \item Derive the wave equations for electric and magnetic fields, emphasizing the propagating nature of electromagnetic waves.
    \end{itemize}
\end{enumerate}

\subsection*{V. Plane Waves in Dielectric Media}

\subsubsection*{Statements:}

\begin{enumerate}
    \item \textbf{Propagation of Electromagnetic Waves in Dielectric Media:}
    \begin{itemize}
        \item Explore the behavior of electromagnetic waves, particularly plane waves, in dielectric media.
    \end{itemize}
\end{enumerate}

\subsection*{VI. Poynting Theorem and Poynting Vector}

\subsubsection*{Statements:}

\begin{enumerate}
    \item \textbf{Introduction to Poynting Theorem:}
    \begin{itemize}
        \item Introduce Poynting's theorem as a statement of energy conservation in electromagnetic fields.
    \end{itemize}
    
    \item \textbf{Poynting Vector:}
    \begin{itemize}
        \item Define the Poynting vector as the directional energy flux per unit area associated with an electromagnetic field.
    \end{itemize}
\end{enumerate}

\section*{Note to Teachers:}

\begin{itemize}
    \item Use visual aids, such as diagrams and animations, to illustrate the concepts of Maxwell's equations and their derivations.
    \item Conduct experiments or demonstrations to show the physical manifestations of electromagnetic waves.
    \item Connect Maxwell's equations to technological applications, emphasizing their role in modern communications.
    \item Assign problems or examples related to the application of Maxwell's equations to reinforce understanding.
\end{itemize}


\rule{\linewidth}{0.5mm}

\subsection*{I. Plane Electromagnetic Waves Through Vacuum}

\subsubsection*{Statements:}

\begin{enumerate}
    \item \textbf{Introduction to Plane Electromagnetic Waves:}
    \begin{itemize}
        \item Introduce the concept of plane electromagnetic waves and their propagation through a vacuum.
    \end{itemize}
    
    \item \textbf{Transverse Nature of Plane Electromagnetic Waves:}
    \begin{itemize}
        \item Emphasize the transverse nature of electromagnetic waves and the orthogonal relationship between electric and magnetic fields.
    \end{itemize}
\end{enumerate}

\subsection*{II. Plane Electromagnetic Waves Through Isotropic Dielectric Medium}

\subsubsection*{Statements:}

\begin{enumerate}
    \item \textbf{Propagation Through Dielectric Medium:}
    \begin{itemize}
        \item Discuss the behavior of plane electromagnetic waves as they pass through an isotropic dielectric medium.
    \end{itemize}
    
    \item \textbf{Refractive Index and Dielectric Constant:}
    \begin{itemize}
        \item Define refractive index and dielectric constant in the context of wave propagation through dielectric media.
    \end{itemize}
    
    \item \textbf{Wave Impedance:}
    \begin{itemize}
        \item Introduce wave impedance as a characteristic property of the medium affecting the propagation of electromagnetic waves.
    \end{itemize}
\end{enumerate}

\subsection*{III. Propagation Through Conducting Media}

\subsubsection*{Statements:}

\begin{enumerate}
    \item \textbf{Propagation Through Conducting Media:}
    \begin{itemize}
        \item Explore the challenges and characteristics of electromagnetic wave propagation through conducting media.
    \end{itemize}
    
    \item \textbf{Relaxation Time:}
    \begin{itemize}
        \item Define relaxation time as a parameter describing the response of conducting media to varying electric fields.
    \end{itemize}
    
    \item \textbf{Skin Depth:}
    \begin{itemize}
        \item Introduce the concept of skin depth as the depth at which the amplitude of electromagnetic waves decreases in a conducting medium.
    \end{itemize}
\end{enumerate}

\section*{Note to Teachers:}

\begin{itemize}
    \item Utilize visual aids, such as animations or simulations, to demonstrate the behavior of electromagnetic waves in different media.
    \item Conduct experiments or demonstrations to show the impact of different media on wave propagation.
    \item Relate the concepts to real-world applications, such as communication systems and material characterization.
    \item Assign problems or examples related to wave propagation in different media to reinforce understanding.
\end{itemize}

\rule{\linewidth}{0.5mm}

\subsection*{I. Boundary Conditions at a Plane Interface Between Two Media}

\subsubsection*{Statements:}

\begin{enumerate}
    \item \textbf{Introduction to Boundary Conditions:}
    \begin{itemize}
        \item Explain the importance of establishing boundary conditions at the interface between two media.
    \end{itemize}
    
    \item \textbf{Boundary Conditions for Electric Field:}
    \begin{itemize}
        \item Define the boundary conditions for the electric field at a plane interface between two dielectric media.
    \end{itemize}
    
    \item \textbf{Boundary Conditions for Magnetic Field:}
    \begin{itemize}
        \item Define the boundary conditions for the magnetic field at a plane interface between two dielectric media.
    \end{itemize}
\end{enumerate}

\subsection*{II. Reflection and Refraction of Plane Waves at a Plane Interface}

\subsubsection*{Statements:}

\begin{enumerate}
    \item \textbf{Laws of Reflection and Refraction:}
    \begin{itemize}
        \item State the laws governing the reflection and refraction of plane waves at a plane interface between two dielectric media.
    \end{itemize}
    
    \item \textbf{Fresnel's Formulae:}
    \begin{itemize}
        \item Introduce Fresnel's formulae as mathematical expressions describing the reflection and transmission coefficients for perpendicular and parallel polarization cases.
    \end{itemize}
    
    \item \textbf{Brewster's Law:}
    \begin{itemize}
        \item Present Brewster's law as a condition for the absence of reflection for a specific angle of incidence.
    \end{itemize}
\end{enumerate}

\section*{Note to Teachers:}

\begin{itemize}
    \item Use diagrams or animations to illustrate the reflection and refraction of plane waves at a dielectric interface.
    \item Conduct experiments or demonstrations to visualize Brewster's law and Fresnel's formulae.
    \item Relate the concepts to real-world applications, such as optics and communication systems.
    \item Assign problems or examples involving reflection and refraction at interfaces to reinforce understanding.
\end{itemize}

\rule{\linewidth}{0.5mm}

\subsection*{I. Description of Polarization}

\subsubsection*{Statements:}

\begin{enumerate}
    \item \textbf{Introduction to Polarization:}
    \begin{itemize}
        \item Define polarization as the orientation of the electric field vector in an electromagnetic wave.
    \end{itemize}
    
    \item \textbf{Linear Polarization:}
    \begin{itemize}
        \item Describe linear polarization, where the electric field oscillates in a straight line.
    \end{itemize}
    
    \item \textbf{Circular Polarization:}
    \begin{itemize}
        \item Describe circular polarization, where the electric field vector rotates in a circular motion.
    \end{itemize}
    
    \item \textbf{Elliptical Polarization:}
    \begin{itemize}
        \item Describe elliptical polarization, a combination of linear and circular polarizations resulting in an elliptical trajectory.
    \end{itemize}
\end{enumerate}

\subsection*{II. Uniaxial and Biaxial Crystals}

\subsubsection*{Statements:}

\begin{enumerate}
    \item \textbf{Introduction to Crystal Polarization:}
    \begin{itemize}
        \item Introduce the polarization properties of crystals, particularly uniaxial and biaxial crystals.
    \end{itemize}
    
    \item \textbf{Double Refraction:}
    \begin{itemize}
        \item Explain the phenomenon of double refraction in crystals, where a single incident ray splits into two.
    \end{itemize}
    
    \item \textbf{Polarization by Double Refraction:}
    \begin{itemize}
        \item Describe how double refraction induces polarization in crystals.
    \end{itemize}
    
    \item \textbf{Nicol Prism:}
    \begin{itemize}
        \item Explain the function of a Nicol prism in producing polarized light from unpolarized light.
    \end{itemize}
    
    \item \textbf{Ordinary and Extraordinary Refractive Indices:}
    \begin{itemize}
        \item Define ordinary and extraordinary refractive indices in the context of double-refracting crystals.
    \end{itemize}
\end{enumerate}

\subsection*{III. Production and Detection of Polarized Light}

\subsubsection*{Statements:}

\begin{enumerate}
    \item \textbf{Production of Polarized Light:}
    \begin{itemize}
        \item Describe methods for producing polarized light, including the use of polarizers.
    \end{itemize}
    
    \item \textbf{Detection of Polarized Light:}
    \begin{itemize}
        \item Explain methods for detecting and analyzing polarized light.
    \end{itemize}
\end{enumerate}

\subsection*{IV. Phase Retardation Plates}

\subsubsection*{Statements:}

\begin{enumerate}
    \item \textbf{Introduction to Phase Retardation Plates:}
    \begin{itemize}
        \item Introduce phase retardation plates as devices that delay one component of a polarized wave with respect to the other.
    \end{itemize}
    
    \item \textbf{Quarter-Wave Plates:}
    \begin{itemize}
        \item Define quarter-wave plates and explain their role in altering the polarization state of light.
    \end{itemize}
    
    \item \textbf{Half-Wave Plates:}
    \begin{itemize}
        \item Define half-wave plates and explain their function in changing the polarization of light.
    \end{itemize}
\end{enumerate}

\subsection*{V. Babinet Compensator and its Uses}

\subsubsection*{Statements:}

\begin{enumerate}
    \item \textbf{Introduction to Babinet Compensator:}
    \begin{itemize}
        \item Introduce the Babinet compensator as a device for compensating phase differences in polarized light.
    \end{itemize}
    
    \item \textbf{Uses of Babinet Compensator:}
    \begin{itemize}
        \item Explain the applications and uses of the Babinet compensator in optical measurements.
    \end{itemize}
\end{enumerate}

\subsection*{VI. Analysis of Polarized Light}

\subsubsection*{Statements:}

\begin{enumerate}
    \item \textbf{Analyzing Polarized Light:}
    \begin{itemize}
        \item Describe methods and devices for analyzing the properties of polarized light.
    \end{itemize}
\end{enumerate}

\section*{Note to Teachers:}

\begin{itemize}
    \item Use visual aids, such as diagrams and animations, to illustrate the different types of polarization and crystal properties.
    \item Conduct hands-on experiments or demonstrations to show the effects of phase retardation plates and Babinet compensators.
    \item Relate the concepts to real-world applications, such as optics and material analysis.
    \item Assign problems or examples involving the manipulation and analysis of polarized light to reinforce understanding.
\end{itemize}

\rule{\linewidth}{0.5mm}

\subsection*{I. Optical Rotation}

\subsubsection*{Statements:}

\begin{enumerate}
    \item \textbf{Introduction to Optical Rotation:}
    \begin{itemize}
        \item Define optical rotation as the phenomenon where the plane of polarized light rotates as it passes through certain substances.
    \end{itemize}
    
    \item \textbf{Biot’s Laws for Rotatory Polarization:}
    \begin{itemize}
        \item Present Biot’s laws, which describe the relationship between the substance, the path length, and the angle of rotation.
    \end{itemize}
\end{enumerate}

\subsection*{II. Fresnel’s Theory of Optical Rotation}

\subsubsection*{Statements:}

\begin{enumerate}
    \item \textbf{Fresnel’s Theory Overview:}
    \begin{itemize}
        \item Introduce Fresnel’s theory, explaining his conceptualization of optical rotation.
    \end{itemize}
    
    \item \textbf{Calculation of Angle of Rotation:}
    \begin{itemize}
        \item Explain the mathematical approach Fresnel used to calculate the angle of rotation based on the properties of the substance.
    \end{itemize}
    
    \item \textbf{Experimental Verification of Fresnel’s Theory:}
    \begin{itemize}
        \item Discuss experiments and methods used to verify Fresnel’s theory in a laboratory setting.
    \end{itemize}
    
    \item \textbf{Specific Rotation:}
    \begin{itemize}
        \item Define specific rotation as the angle of rotation per unit length and concentration of the substance.
    \end{itemize}
\end{enumerate}

\section*{Note to Teachers:}

\begin{itemize}
    \item Use visual aids, such as diagrams and animations, to illustrate optical rotation and the principles of Fresnel’s theory.
    \item Conduct hands-on experiments or demonstrations to show optical rotation and verify Fresnel’s theory.
    \item Relate the concepts to real-world applications, such as pharmaceutical and chemical industries.
    \item Assign problems or examples involving the calculation of optical rotation and verification of Fresnel’s theory to reinforce understanding.
\end{itemize}







%------------Body-Ends--------------------
%\bibliography{ref.bib}
%\bibliographystyle{IEEEtran}
\end{document}
